\section{Related Work}
\label{sec:related}

\subsection{Static Model Merging}
Model merging has emerged as a prominent direction for resource-efficient multi-task learning~\cite{Wortsman2022}. Early works demonstrated that simple weight averaging (or Task Arithmetic) can combine task vectors to edit pre-trained models with zero training cost~\cite{Ilharco2022}. However, as the number of tasks grows, static merging suffers from severe \emph{catastrophic interference} (or task conflict), where updating weights for one task degrades performance on another. To address parameter conflict, TIES-Merging~\cite{Yadav2023} filters low-magnitude parameter changes, resolves sign disagreement across task vectors, and averages only consolidated weights. Other static methods incorporate data-free structural properties such as Fisher information scaling~\cite{Matena2022}, activation-covariance alignment (RegMean)~\cite{jin2022dataless}, or joint pruning and merging (ZipMerge)~\cite{ZipMerge2025}. While these static paradigms are valuable, they generate a singular, immutable parameter set for all inputs, fundamentally capping multi-task capacity.

\subsection{Dynamic Model Merging and Routing}
To surpass the capacity ceiling of static merging, dynamic model merging computes sample-specific blending coefficients at runtime. This paradigm draws inspiration from Mixture of Experts (MoE) architectures~\cite{shazeer2017outrageously, fedus2022switch}, which route tokens or entire samples to distinct feed-forward subnetworks. However, standard MoE requires training massive networks from scratch, which is computationally prohibitive. 

In the post-training merge setting, AdaMerging~\cite{AdaMerging2023} introduced adaptive test-time model merging by optimizing blending coefficients on-the-fly using unsupervised entropy minimization. While innovative, test-time adaptation (TTA) can be slow at inference and highly sensitive to batch sizes and noisy test streams~\cite{liao2021are, OFS2025}. To circumvent test-time latency, QWS-Merge~\cite{vance2025quantum} proposed training a gating network during a quick offline calibration phase. However, as discussed in Section~\ref{sec:intro}, QWS-Merge introduced convoluted, quantum-inspired mechanisms (such as task eigenstates and wave projection operators) under the premise that classical linear routers collapse on out-of-distribution tasks like SVHN. In contrast, our work demonstrates that this collapse is not a structural failure of linear networks but a standard, preventable overfitting phenomenon, which can be elegantly resolved via classical regularization.

\subsection{Regularization in Parameter Spaces}
Regularization is the cornerstone of machine learning, serving to constrain hypothesis spaces and improve out-of-distribution generalization~\cite{goodfellow2016deep, shwartz2014understanding}. Timeless techniques such as $L_2$ weight regularization (weight decay)~\cite{Loshchilov2019} and softmax temperature scaling~\cite{srivastava2014dropout} are ubiquitous in deep neural network training. In the context of model merging, previous works have applied regularization to the underlying model landscapes (e.g., RegCalMerge~\cite{RegCalMerge2024} and SAIM~\cite{saim}) or parameterized trajectories (e.g., PolyMerge~\cite{PolyMerge2025}). However, these works applied regularization to the high-dimensional parameter spaces of the models themselves. Surprisingly, while gating network regularizations (such as load-balancing losses, entropy-regularized gating, routing noise, and gating dropout) are foundational pillars in traditional sparse Mixture-of-Experts (MoE) architectures to prevent expert collapse and stabilize routing~\cite{shazeer2017outrageously, fedus2022switch}, regularization of the \emph{routing network itself} has been almost entirely overlooked in post-hoc dynamic parameter fusion. Our work fills this gap by adapting these classic MoE gating regularization principles directly to the post-hoc merging paradigm. Rather than designing complex routing objectives, we demonstrate that simple, classic $L_2$ weight regularization and temperature-scaled softmax scaling applied to a 768-parameter linear router are sufficient to unlock highly robust and competitive performance under distribution shifts, rendering complex dynamic fusion architectures obsolete.
